% LSPV-2 REPORT::DELAYOR1::2026-10-04T15:52:00-07:00
% !TEX program = xelatex
\documentclass[11pt]{article}
\usepackage[margin=.86in]{geometry}
\usepackage{amsmath,amssymb}
\begin{document}

\section*{LSPV-2 DELAYOR1: DELAY-wrapper coefficient barrier and failure of the direct Williams OR transplant}

\textbf{GLOBAL STATUS: OPEN.} P versus NP remains open. This report continues RETURNSQUEEZE1 by attacking the fine-grained return bridge. It proves two restricted no-go lemmas and one exact exponent calculation. It does not prove that all nonblackbox return bridges fail.

\section*{CORPUS REVIEWED}
The active basis is verified RETURNSQUEEZE1 together with ENVELOPE1, TOTALENV1, QLEN1, QUOTIENT1, EXCESS1, TURANEX1, PAIRDEG1, VIEWSEP1, TASK-001, TASK-002, TASK-004, TASK-005, root main.tex, and MAIN0901. The current Boundary-Squeeze target is a repeatable SAT-to-reconstruction-to-SAT transformation with
\[
K'\le C_T+\theta K,\qquad \theta<1,
\]
and a compatible space contraction.

Primary literature checked for this run includes Williams' STOC 2024 self-improvement theorem, Santhanam--Srinivasan on limits of CNF sparsification, Fortnow--Santhanam and Drucker on instance/OR compression, Dell--van Melkebeek on SAT sparsification barriers, Buss--Williams on alternation-trading limits, and classical program-speedup/program-analysis literature for the wrapper audit.

\section*{Literature lesson}
Williams' successful Circuit-SAT self-improvement uses two properties that are not supplied by ordinary CNF self-reduction:
\begin{enumerate}
\item the assignment space is embarrassingly parallel and can be packed by an OR construction; and
\item the solving algorithm can be modeled extremely efficiently inside an instance of the same circuit-analysis problem.
\end{enumerate}
Williams explicitly states that the construction is highly nonblackbox: descriptions of circuits solving Circuit-SAT are embedded into other Circuit-SAT instances. The paper also states as an open problem whether self-improvement can reach the P=NP regime.

Santhanam--Srinivasan prove unconditional representation lower bounds showing that some CNFs cannot be written as subexponential ORs of much smaller CNFs. Their theorem concerns equality of represented Boolean functions, not merely preservation of satisfiability, so it is a hostile representation result rather than a direct lower bound on a SAT return reduction.

Fortnow--Santhanam, Drucker, and Dell--van Melkebeek give strong conditional barriers to general instance/OR compression and sparsification. These cannot be terminal contradictions under the counterfactual P=NP, but they confirm that the required compression is a strong operation.

\section*{PROVED/EXACT: DELAYCOEF1 blackbox coefficient barrier}
Assume there exists one exact deterministic SAT decider \(B\) running in polynomial time. For every sufficiently large integer \(L\), define an extensionally equivalent decider \(A_L\) that runs \(B\), stores the correct answer, and then performs dummy work so that on length-\(m\) inputs its running time is
\[
\Theta(m^L).
\]
Thus \(A_L\) is a correct exact SAT decider for every \(L\), differing from \(B\) only by a DELAY wrapper.

Consider a deterministic oracle-style transformation \(T\) that is extensional in its use of the SAT decider: its oracle queries depend on the input and previous oracle answer bits, but not on the source code or timing behavior of the oracle implementation. Suppose there is a constant \(\delta>0\) and infinitely many input lengths \(n\) for which the execution of \(T^{A_L}\) on some length-\(n\) input makes an actually executed oracle call of length
\[
n^{\delta-o(1)}.
\]
Because every \(A_L\) returns the same SAT answers, the extensional query transcript is the same for all \(L\).

That call alone takes
\[
n^{\delta L-o(L)}
\]
steps under \(A_L\). Therefore, if a claimed uniform exponent theorem has
\[
T^{A_L}(n)\le n^{C+\theta L+o(1)}
\]
with constants \(C,\theta\) independent of \(L\), then for arbitrarily large \(L\),
\[
\delta L\le C+\theta L+o(L).
\]
Dividing by \(L\) and letting \(L\to\infty\) gives
\[
\boxed{\theta\ge\delta.}
\]

Hence a wrapper-robust blackbox exponent contraction with \(\theta<1\) cannot make any near-original-scale call of length \(n^{1-o(1)}\). More generally, its inherited coefficient is at least the largest power-scale exponent of any actually executed blackbox call.

This is a theorem about extensional blackbox transformations. A nonblackbox transformation that inspects code and changes behavior on DELAY-wrapped implementations is not covered.

\section*{Connection to RETURNSQUEEZE1}
RETURNSQUEEZE1 gives the direct-route upper-bound coefficient
\[
\theta\le\lambda b
\]
when the return bridge maps length \(n\) to reconstruction length
\[
N=n^{\lambda+o(1)}
\]
and inherited SAT queries have maximum actual scale
\[
N^{b+o(1)}.
\]
If this maximum is attained on the relevant worst-case executions and the bridge uses the SAT solver extensionally, DELAYCOEF1 gives the matching necessity
\[
\theta\ge\lambda b.
\]
Thus in that blackbox fixed-parameter regime, the coefficient \(\lambda b\) is not merely accounting slack: it is wrapper-tight.

The same argument applies to any evaluator route whose actual call-size exponent can be source-certified.

\section*{PROVED/EXACT: RESTRICTCOVER1}
Let
\[
\rho_1,\ldots,\rho_q
\]
be partial assignments on \(n\) Boolean variables whose subcubes cover the entire Boolean cube:
\[
\bigcup_{i=1}^q[\rho_i]=\{0,1\}^n.
\]
If every \(\rho_i\) leaves at most \(m\) variables free, then each subcube contains at most \(2^m\) assignments. Therefore
\[
2^n\le q2^m,
\]
and hence
\[
\boxed{m\ge n-\log_2q.}
\]
In particular, if \(q\le n^c\) for a constant \(c\), then
\[
\boxed{m\ge n-c\log_2n.}
\]
Thus polynomially many explicit restriction branches cannot reduce the assignment dimension from \(n\) to \(n^\gamma\) for any fixed \(\gamma<1\) while covering all assignments. A Williams-style split that reduces the remaining variable dimension by a constant fraction necessarily represents exponentially many branches somewhere.

RESTRICTCOVER1 is a subcube-cover theorem only. Formula-dependent semantic pruning may avoid a full cube cover; no arbitrary SAT lower bound is inferred.

\section*{PROVED/EXACT exponent calculation: direct Williams OR transplant to ordinary SAT}
Take a SAT representation of explicit length \(N\) and split off \(h\) Boolean choices. The literal Williams OR construction packs all \(2^h\) branches into a circuit of size
\[
S'=2^{h}N^{1+o(1)}
\]
up to model-dependent low-order factors.

If this packed circuit is returned to an ordinary SAT solver by a linear or near-linear Tseitin/Cook--Levin style encoding, the resulting SAT instance has length
\[
M=S'^{1+o(1)}.
\]
An assumed ordinary SAT decider of exponent \(K\) then takes
\[
M^{K+o(1)}
=
2^{hK}N^{K+o(1)}.
\]
Compared with simply running the same decider on the original length-\(N\) instance, this packing step does not decrease the polynomial exponent and incurs an additional factor \(2^{hK}\).

Therefore
\[
\boxed{\text{the literal Williams OR packing plus generic return-to-CNF is not the missing polynomial-exponent contraction.}}
\]
The reason Williams' theorem succeeds for Circuit-SAT is different: the assumed algorithm is applied directly inside the circuit-analysis representation, and the previous-stage solver can be modeled efficiently in that same representation. Ordinary CNF return via generic computation encoding loses precisely that advantage.

Status: \textbf{FALSIFIED/RETIRED} for the literal transplant described above. This does not rule out a source-specific nonblackbox packing that remains compact in the exact input language consumed by the improved solver.

\section*{Representation evidence from CNF sparsification}
Santhanam--Srinivasan prove that for every superlinear \(f(n)\) there are \(n\)-variable CNFs of size at most \(f(n)\) that cannot be represented as the OR of \(2^{n-\varepsilon n}\) linear-size CNFs for any fixed \(\varepsilon>0\); they also prove a hierarchy of non-sparsifiability. This shows that an unrestricted functional OR-packing theorem for CNFs is false.

However, their lower bound preserves the entire Boolean function. A SAT return bridge needs only preserve whether the function has at least one satisfying assignment. Under P=NP, semantic answer-preserving compression to a constant is trivial after paying the original decision cost. Therefore the representation theorem cannot replace the missing charged-construction theorem.

\section*{Hostile controls and barrier audit}
\textbf{DELAY.} DELAYCOEF1 is designed specifically to prevent a universal blackbox contraction from silently assuming a fastest implementation of the hypothetical SAT decider.

\textbf{DEAD/COPY/padding.} These can alter nominal instance-size exponents. The product of the outer size scale and actual call scale, not a padded local exponent, must be charged.

\textbf{INTERP/RECOMP.} Nonblackbox code transformation may try to strip delays or choose an alternate implementation. DELAYCOEF1 then no longer applies directly. A surviving theorem must specify the exact allowed source inspection and prove its correctness and cost. General program-analysis/Rice-style impossibility results warn against assuming a complete semantic optimizer, but they do not by themselves rule out a SAT-specific source transformation.

\textbf{Semantic oracle.} A semantic SAT oracle can compress to a one-bit answer immediately. This demonstrates again that semantic existence of a short equivalent object is not the acquisition theorem.

\textbf{Algebraic oracle.} The DELAY and subcube-cover arguments survive algebraic information as accounting/cover statements, but neither supplies the final nonalgebrizing hinge.

\textbf{Easy families.} Horn, 2-SAT, affine/Tseitin, bounded-treewidth, and independent-product formulas may admit compact semantic projection. Therefore compact projection itself is not a hardness resource.

\section*{What changed}
The Boundary-Squeeze branch is narrowed to a nonblackbox return theorem. A fixed-parameter extensional blackbox reduction cannot beat the largest power-scale SAT query it actually makes, because DELAY wrappers make that coefficient unavoidable. Polynomially many explicit restriction branches also cannot shrink assignment dimension by a power. The direct Williams OR packing becomes larger when generically returned to ordinary CNF and therefore does not contract an unknown polynomial exponent.

\section*{ADOPTED / REJECTED / RETIRED}
\textbf{Adopted:} RETURNSQUEEZE1, Williams self-improvement as the closest positive architecture, RESTRICTCOVER1, DELAYCOEF1.

\textbf{Rejected as sufficient:} ordinary restriction branching, polynomially many explicit subcubes, generic OR packing followed by Tseitin/Cook--Levin return, function-representation lower bounds as decision-runtime lower bounds.

\textbf{OPEN rescue:} a source-specific nonblackbox transformation which can inspect the solver or the exact SU/vMS reconstruction, avoid original-scale blackbox calls, pack exponentially many semantic branches implicitly, and emit an ordinary SAT instance whose actual written length gives a composite coefficient below one.

\section*{Single next target}
The next theorem should be stated nonblackbox from the outset. Let \(\mathcal T\) receive source code for an arbitrary exact SAT decider \(A\) with a certified polynomial clock. Seek either:
\[
\boxed{
\mathcal T(A)\text{ exact SAT},\qquad
K_{\mathcal T(A)}\le C+\theta K_A,\quad\theta<1
}
\]
with explicit invariance under DELAY/INTERP/RECOMP, or a diagonal/wrapper obstruction proving that no source transformation in a precisely declared compiler class can achieve this.

The most promising positive analogy is Williams' same-representation self-modeling. The corresponding LSPV question is whether the exact reconstruction/product-span structure gives a representation in which a prior-stage SAT solver can be modeled without generic Cook--Levin expansion and without paying the solver exponent again. That remains OPEN.

\section*{References}
R. Ryan Williams, Self-Improvement for Circuit-Analysis Problems, STOC 2024.

R. Santhanam and S. Srinivasan, On the Limits of Sparsification, ECCC TR11-131, 2011.

L. Fortnow and R. Santhanam, Infeasibility of Instance Compression and Succinct PCPs for NP, JCSS 2011.

A. Drucker, New Limits to Classical and Quantum Instance Compression, SIAM J. Comput. 2015.

H. Dell and D. van Melkebeek, Satisfiability Allows No Nontrivial Sparsification Unless the Polynomial-Time Hierarchy Collapses, JACM 2014.

S. Buss and R. Ryan Williams, Limits on Alternation-Trading Proofs for Time-Space Lower Bounds, Computational Complexity 2015.

\section*{FINAL STATUS}
DELAYCOEF1 and RESTRICTCOVER1 are PROVED/EXACT in their stated restricted models. The direct Williams-OR-to-CNF transplant is FALSIFIED/RETIRED as the missing contraction. A nonblackbox source-specific return bridge, source-certified query exponents, repeatability, time-space fixed-point collision, arbitrary-algorithm obligation, nonrelativization, nonalgebrization, and P versus NP all remain OPEN.

% END_OF_VERIFIED_REPORT::LSPV-2::DELAYOR1
\end{document}
